% =====================================================================
\chapter{Derivadas matriciales: capa lineal y softmax}
\label{tema:3}
% =====================================================================

Casi todo el cómputo del entrenamiento de un modelo de lenguaje consiste en derivar dos tipos de operación: capas lineales y la combinación de softmax con entropía cruzada en la salida. Este tema define el gradiente de una función escalar respecto de una matriz, obtiene los gradientes de la capa lineal, individualmente y por lotes, y deduce el resultado más usado del curso: el gradiente de la entropía cruzada respecto de los logits es $p-y$. Se sigue~\cite[§5.3--§5.5]{mml}, la tarea~2 de~\cite{cs224n} y, como tabla de consulta, \cite{cookbook}.

\section{Gradiente respecto de una matriz}

\begin{definicion}[Gradiente respecto de una matriz]\label{def:grad-matriz}
Sea $L\colon\R^{m\times n}\to\R$. Su \emph{gradiente} en $W$ es la matriz $\nabla_W L=\partial L/\partial W\in\R^{m\times n}$ de entradas
\[
\Bigl(\frac{\partial L}{\partial W}\Bigr)_{ij}=\frac{\partial L}{\partial W_{ij}}.
\]
Equivale a identificar $\R^{m\times n}$ con $\R^{mn}$ y reordenar el gradiente con la forma de $W$.
\end{definicion}

\begin{proposicion}[Aproximación lineal]\label{prop:frobenius}
Si $L$ es diferenciable en $W$, para toda perturbación $E\in\R^{m\times n}$,
\[
L(W+E)=L(W)+\Bigl\langle\frac{\partial L}{\partial W},E\Bigr\rangle+o(\norm{E}),\qquad
\langle A,B\rangle=\sum_{i,j}A_{ij}B_{ij}=\tr(A\T B).
\]
\end{proposicion}

\begin{proof}
Es la \cref{def:diferenciable} en $\R^{mn}$: el término lineal $\nabla L\T h$ es $\sum_{i,j}(\partial L/\partial W_{ij})E_{ij}$. La igualdad $\sum_{i,j}A_{ij}B_{ij}=\tr(A\T B)$ se obtiene desarrollando $(A\T B)_{jj}=\sum_i A_{ij}B_{ij}$ y sumando en $j$.
\end{proof}

La regla práctica que se deriva es que \emph{el gradiente de un escalar respecto de un parámetro tiene la misma forma que el parámetro}. Comprobar las formas es la primera verificación de cualquier fórmula de este tema.

\section{La capa lineal}

\begin{proposicion}[Gradientes de una capa lineal]\label{prop:capa}
Sean $W\in\R^{m\times n}$, $x\in\R^{n}$, $b\in\R^{m}$, $y=Wx+b$ y $L=\ell(y)$ con $\ell\colon\R^{m}\to\R$ diferenciable. Sea $\delta=\nabla_y\ell\in\R^{m}$ la \emph{señal de error} que llega a la salida. Entonces
\[
\frac{\partial L}{\partial W}=\delta\,x\T,\qquad
\frac{\partial L}{\partial b}=\delta,\qquad
\frac{\partial L}{\partial x}=W\T\delta.
\]
\end{proposicion}

\begin{proof}
Como $y_k=\sum_{j}W_{kj}x_j+b_k$, se tiene $\partial y_k/\partial W_{ij}=\mathbb{I}[k=i]\,x_j$: la entrada $W_{ij}$ solo afecta a la salida $y_i$. Por el \cref{cor:cadena-escalar},
\[
\frac{\partial L}{\partial W_{ij}}=\sum_{k}\delta_k\,\frac{\partial y_k}{\partial W_{ij}}=\delta_i\,x_j=(\delta x\T)_{ij}.
\]
Del mismo modo, $\partial y_k/\partial b_i=\mathbb{I}[k=i]$ da $\partial L/\partial b_i=\delta_i$. Por último, $J_y(x)=W$ (\cref{ej:jac-lineal}) y el \cref{cor:cadena-escalar} da $\partial L/\partial x=W\T\delta$.
\end{proof}

Las formas encajan: $\delta x\T$ es $(m\times 1)(1\times n)=m\times n$, como $W$, y $W\T\delta$ es $(n\times m)(m\times 1)=n\times 1$, como $x$. El vector $W\T\delta$ es la señal de error que la capa transmite a la capa anterior.

\begin{ejemplo}\label{ej:capa}
Sean $W=\begin{pmatrix}1&2\\0&1\end{pmatrix}$, $x=(1,-1)\T$, $b=0$, un objetivo $t=(0,1)\T$ y $L=\tfrac12\norm{y-t}^{2}$, de modo que $\delta=y-t$.
\begin{center}\small
\begin{tabular}{@{}lll@{}}
\toprule
Paso & Cálculo & Resultado \\
\midrule
Hacia delante & $y=Wx$ & $(-1,-1)\T$ \\
Pérdida & $\tfrac12\bigl((-1)^{2}+(-2)^{2}\bigr)$ & $2{,}5$ \\
Señal de error & $\delta=y-t$ & $(-1,-2)\T$ \\
Gradiente de $W$ & $\delta x\T$ & $\begin{pmatrix}-1&1\\-2&2\end{pmatrix}$ \\[2mm]
Gradiente de $x$ & $W\T\delta$ & $(-1,-4)\T$ \\
\bottomrule
\end{tabular}
\end{center}
\end{ejemplo}

\begin{corolario}[Capa lineal por lotes]\label{cor:lotes}
Sean $X\in\R^{B\times n}$ una matriz cuyas filas son $B$ ejemplos, $W\in\R^{n\times m}$, $Z=XW\in\R^{B\times m}$ y $L=\ell(Z)$ con $\Delta=\partial L/\partial Z\in\R^{B\times m}$. Entonces
\[
\frac{\partial L}{\partial W}=X\T\Delta,\qquad\frac{\partial L}{\partial X}=\Delta\,W\T.
\]
\end{corolario}

\begin{proof}
Como $Z_{bk}=\sum_j X_{bj}W_{jk}$, se tiene $\partial Z_{bk}/\partial W_{jl}=\mathbb{I}[k=l]\,X_{bj}$, y por la regla de la cadena $\partial L/\partial W_{jl}=\sum_{b}\Delta_{bl}X_{bj}=(X\T\Delta)_{jl}$. Análogamente, $\partial Z_{bk}/\partial X_{cj}=\mathbb{I}[b=c]\,W_{jk}$ da $\partial L/\partial X_{cj}=\sum_k\Delta_{ck}W_{jk}=(\Delta W\T)_{cj}$.
\end{proof}

El \cref{cor:lotes} es la \cref{prop:capa} escrita con filas en lugar de columnas y sumada sobre los ejemplos del lote: $X\T\Delta=\sum_b x_b\,\delta_b\T$. Es la convención de las bibliotecas de aprendizaje profundo.

\begin{observacion}[Conexión con los LLM]
Las proyecciones de consultas, claves y valores de la atención (módulo~13), las dos capas del bloque MLP y la capa final que produce los logits son capas lineales. El paso hacia atrás de un Transformer se compone, en su mayor parte, de los productos $X\T\Delta$ y $\Delta W\T$.
\end{observacion}

\section{El jacobiano de la softmax}

Se recuerda del módulo~1 que, para $z\in\R^{K}$, $p=\softmax(z)$ tiene componentes $p_i=e^{z_i}/S$ con $S=\sum_k e^{z_k}$.

\begin{proposicion}[Jacobiano de la softmax]\label{prop:jac-softmax}
\[
\frac{\partial p_i}{\partial z_j}=p_i\bigl(\mathbb{I}[i=j]-p_j\bigr),\qquad\text{es decir,}\qquad J_{\softmax}(z)=\diag(p)-pp\T.
\]
Esta matriz es simétrica y cumple $J_{\softmax}(z)\,\uno=0$.
\end{proposicion}

\begin{proof}
Por la regla del cociente, con $\partial S/\partial z_j=e^{z_j}$: si $i=j$, $\partial p_i/\partial z_i=(e^{z_i}S-e^{z_i}e^{z_i})/S^{2}=p_i-p_i^{2}$; si $i\ne j$, $\partial p_i/\partial z_j=(0-e^{z_i}e^{z_j})/S^{2}=-p_ip_j$. La simetría es evidente en $\diag(p)-pp\T$. Por último, $\diag(p)\uno=p$ y $pp\T\uno=p\,(p\T\uno)=p$, porque las componentes de $p$ suman~1.
\end{proof}

La igualdad $J_{\softmax}\uno=0$ es la versión infinitesimal de la invariancia por traslación de la softmax (módulo~1, tema~4): sumar la misma constante a todos los logits no cambia las probabilidades.

\section{Entropía cruzada y el gradiente \texorpdfstring{$p-y$}{p − y}}

Sea $c$ el índice del token correcto e $y=e_c$ su vector \emph{one-hot}. La pérdida de entropía cruzada es $L(z)=-\ln p_c$.

\begin{teorema}[Gradiente de la entropía cruzada]\label{teo:p-y}
\[
\nabla_z L=p-y.
\]
\end{teorema}

\begin{proof}
\emph{Primera demostración (regla de la cadena).} $L$ depende de $z$ solo a través de $p_c$, con $\partial L/\partial p_c=-1/p_c$. Por la \cref{prop:jac-softmax},
\[
\frac{\partial L}{\partial z_j}=-\frac{1}{p_c}\,\frac{\partial p_c}{\partial z_j}=-\frac{1}{p_c}\,p_c\bigl(\mathbb{I}[c=j]-p_j\bigr)=p_j-y_j.
\]
\emph{Segunda demostración (log-suma-exp).} $L=-\ln\bigl(e^{z_c}/S\bigr)=-z_c+\LSE(z)$. La derivada de $-z_c$ respecto de $z_j$ es $-y_j$, y la de $\LSE(z)=\ln S$ es $e^{z_j}/S=p_j$ (módulo~1, tema~4).
\end{proof}

\begin{ejemplo}\label{ej:p-y}
Sean $z=(2,1,0)\T$ y $c=1$.
\begin{center}\small
\begin{tabular}{@{}lccc@{}}
\toprule
 & Token 1 (correcto) & Token 2 & Token 3 \\
\midrule
$p=\softmax(z)$ & $0{,}665$ & $0{,}245$ & $0{,}090$ \\
$y$ & $1$ & $0$ & $0$ \\
$\nabla_z L=p-y$ & $-0{,}335$ & $0{,}245$ & $0{,}090$ \\
\bottomrule
\end{tabular}
\end{center}
La pérdida es $L=-\ln 0{,}6652\approx 0{,}408$ nats. Las componentes del gradiente suman~0.
\end{ejemplo}

\begin{observacion}[Interpretación]
Un paso de descenso $z\leftarrow z-\eta(p-y)$ aumenta el logit correcto en $\eta(1-p_c)$ y disminuye cada logit incorrecto en $\eta p_i$: se penaliza más a los competidores más probables. El gradiente es pequeño cuando el modelo acierta con confianza ($p_c\approx 1$) y grande cuando se equivoca. Como la expresión $p-y$ no contiene cocientes por $p_c$, las bibliotecas implementan la softmax y la entropía cruzada como una sola operación, más estable que componer las dos (módulo~11).
\end{observacion}

\begin{corolario}[Capa de salida de un LLM]\label{cor:salida}
Sean $h\in\R^{d}$ el estado oculto de la última posición, $W\in\R^{K\times d}$ la matriz de salida, $z=Wh$ y $L=-\ln\softmax(z)_c$. Entonces
\[
\frac{\partial L}{\partial W}=(p-y)\,h\T,\qquad\frac{\partial L}{\partial h}=W\T(p-y).
\]
\end{corolario}

\begin{proof}
Basta aplicar la \cref{prop:capa} con $\delta=\nabla_z L=p-y$ (\cref{teo:p-y}).
\end{proof}

Con un vocabulario de $K=10^{5}$ tokens, $W$ tiene $10^{5}$ filas; la fila $i$ recibe la corrección $(p_i-y_i)\,h$, de modo que todas las filas se actualizan en cada paso, no solo la del token correcto.

\begin{lecturas}
\item Deisenroth, Faisal y Ong~\cite{mml}, §5.3, §5.4 «Gradients of Matrices» y §5.5 «Useful Identities for Computing Gradients».
\item Stanford CS224N~\cite{cs224n}, tarea~2 (word2vec), apartados escritos: gradientes de la softmax.
\item Petersen y Pedersen~\cite{cookbook}, cap.~2 «Derivatives», como tabla de identidades.
\end{lecturas}

\section*{Ejercicios del tema 3}
\addcontentsline{toc}{section}{Ejercicios del tema 3}

\begin{ejercicio}{1}{eje:3-embeddings}
Un modelo tiene una matriz de \emph{embeddings} $E$ de tamaño $50\,000\times 4096$. ¿Qué forma tiene $\partial L/\partial E$ y cuántos números contiene? ¿Qué filas pueden ser no nulas en un lote concreto?
\end{ejercicio}

\begin{ejercicio}{1}{eje:3-signos}
En el \cref{teo:p-y}, ¿qué signo tiene la componente del token correcto y qué signo tienen las demás? Describa el efecto de un paso de descenso de gradiente sobre cada logit.
\end{ejercicio}

\begin{ejercicio}{2}{eje:3-capa}
Sean $W=\begin{pmatrix}1&0&-1\\2&1&0\end{pmatrix}$, $x=(1,2,1)\T$, $t=(0,3)\T$ y $L=\tfrac12\norm{Wx-t}^{2}$. Calcule $y=Wx$, $L$, $\partial L/\partial W$ y $\partial L/\partial x$.
\end{ejercicio}

\begin{ejercicio}{2}{eje:3-softmax}
Sean $z=(0,\ln 2,\ln 5)\T$ y $c=2$. Calcule $p$, la pérdida $L$ y el gradiente $\nabla_z L$.
\end{ejercicio}

\begin{ejercicio}{3}{eje:3-nucleo}
Demuestre que las componentes de $p-y$ suman cero y que $\uno$ pertenece al núcleo de $J_{\softmax}$. Demuestre además que $\nabla_z L$ es ortogonal a $\uno$ y explique por qué es coherente con la invariancia de la softmax por traslación.
\end{ejercicio}

\begin{ejercicio}{3}{eje:3-word2vec}
En el modelo \emph{skip-gram} de word2vec, la palabra central tiene un vector $v_c\in\R^{d}$ y cada palabra $w$ del vocabulario tiene un vector de contexto $u_w$, columna $w$ de $U\in\R^{d\times\abs{\V}}$. La pérdida para la palabra de contexto observada $o$ es $J=-\ln\hat y_o$, con $\hat y=\softmax(U\T v_c)$. (a) Demuestre que $\partial J/\partial v_c=U(\hat y-y)$, con $y=e_o$. (b) Demuestre que $\partial J/\partial u_w=(\hat y_w-y_w)\,v_c$. (c) Calcule ambos con $d=2$, $\abs{\V}=3$, $u_1=(1,0)\T$, $u_2=(0,1)\T$, $u_3=(1,1)\T$, $v_c=(\ln 2,0)\T$ y $o=2$.
\fuente{adaptado de Stanford CS224N, invierno de 2021~\cite{cs224n}, tarea~2, apartados escritos (b) y (c); el caso numérico es propio.}
\end{ejercicio}
